\documentclass[10pt,a4paper]{amsart}
\usepackage{amsmath,amssymb,amsthm,graphicx}
% Retained preamble for the Stacks Project chapter snapshot in this directory.
% The chapter's own preamble.tex is a copy of the shared Stacks Project preamble; the block
% below is sliced MECHANICALLY, line for line, from it. Nothing was added, removed or
% redefined; the wrapper class, the xr-hyper/hyperref wiring and the externaldocument list
% are omitted because the assembler supplies the document class.
%
% lines taken from preamble.tex: 3, 5, 9, 11, 12, 156, 157, 158, ... (51 lines total)
%
% package coverage restored from the chapter's own preamble (2026-10-13): xy, xy-2cell

\usepackage[all]{xy}
\xyoption{2cell}
\UseAllTwocells

\newenvironment{reference}{\comment}{\endcomment}
\newenvironment{slogan}{\comment}{\endcomment}
\newenvironment{history}{\comment}{\endcomment}
\theoremstyle{plain}
\newtheorem{theorem}[subsection]{Theorem}
\newtheorem{proposition}[subsection]{Proposition}
\newtheorem{lemma}[subsection]{Lemma}
\theoremstyle{definition}
\newtheorem{definition}[subsection]{Definition}
\newtheorem{example}[subsection]{Example}
\newtheorem{exercise}[subsection]{Exercise}
\newtheorem{situation}[subsection]{Situation}
\theoremstyle{remark}
\newtheorem{remark}[subsection]{Remark}
\newtheorem{remarks}[subsection]{Remarks}
\numberwithin{equation}{subsection}
\def\lim{\mathop{\mathrm{lim}}\nolimits}
\def\colim{\mathop{\mathrm{colim}}\nolimits}
\def\Spec{\mathop{\mathrm{Spec}}}
\def\Hom{\mathop{\mathrm{Hom}}\nolimits}
\def\Ext{\mathop{\mathrm{Ext}}\nolimits}
\def\SheafHom{\mathop{\mathcal{H}\!\mathit{om}}\nolimits}
\def\SheafExt{\mathop{\mathcal{E}\!\mathit{xt}}\nolimits}
\def\Sch{\mathit{Sch}}
\def\Mor{\mathop{\mathrm{Mor}}\nolimits}
\def\Ob{\mathop{\mathrm{Ob}}\nolimits}
\def\Sh{\mathop{\mathit{Sh}}\nolimits}
\def\NL{\mathop{N\!L}\nolimits}
\def\CH{\mathop{\mathrm{CH}}\nolimits}
\def\proetale{{pro\text{-}\acute{e}tale}}
\def\etale{{\acute{e}tale}}
\def\QCoh{\mathit{QCoh}}
\def\Ker{\mathop{\mathrm{Ker}}}
\def\Im{\mathop{\mathrm{Im}}}
\def\Coker{\mathop{\mathrm{Coker}}}
\def\Coim{\mathop{\mathrm{Coim}}}
\DeclareMathSymbol{\boxtimes}{\mathbin}{AMSa}{"02}
\def\QCohstack{\mathcal{QC}\!\mathit{oh}}
\def\Cohstack{\mathcal{C}\!\mathit{oh}}
\def\Spacesstack{\mathcal{S}\!\mathit{paces}}
\def\Quotfunctor{\mathrm{Quot}}
\def\Hilbfunctor{\mathrm{Hilb}}
\def\Curvesstack{\mathcal{C}\!\mathit{urves}}
\def\Polarizedstack{\mathcal{P}\!\mathit{olarized}}
\def\Complexesstack{\mathcal{C}\!\mathit{omplexes}}
\def\Pic{\mathop{\mathrm{Pic}}\nolimits}
\def\Picardstack{\mathcal{P}\!\mathit{ic}}
\def\Picardfunctor{\mathrm{Pic}}
\def\Deformationcategory{\mathcal{D}\!\mathit{ef}}
\title{Stacks, Differential Graded Sheaves: lemma-small-acyclics}
\author{The Stacks Project Authors}
\date{stacks-project:sdga.tex:local-snapshot-20261012; GNU FDL 1.2 or later}
\begin{document}
\maketitle
\noindent\textit{Editorial note.} Statement span is lines 3851-3860 of sources/stacks-sdga/pinned-original-source.tex; proof span is lines 3862-3949, the complete author proof (86 body lines).\par\medskip
\noindent\textit{Editorial note.} Navigation stubs: 0 cited-but-not-included labels are supplied as non-load-bearing pointers with their pinned tags where the snapshot tag index has them.\par\medskip
\noindent\textit{Editorial note.} extra\_packages is empty on purpose: the retained preamble already loads xy with options and a second declaration clashes in the assembler wrapper.\par\medskip
\noindent\textit{Editorial note.} Text is retained verbatim, including the author's own typos and placeholders.\par\medskip
\noindent\textit{Editorial note.} Queue index 4 of queue\_sdga.json (pb=3601, stmt 3851-3860, proof 3862-3949).\par\medskip
\noindent\textit{Allowed standard prerequisites (editorial).}\begin{itemize}
\item Explicit prerequisite boundary. The cited inputs are prerequisites; the argument in the retained span is not replaced by any cited result.
\end{itemize}\medskip
\begin{lemma}
\label{lemma-small-acyclics}
Let $(\mathcal{C}, \mathcal{O})$ be a ringed site.
Let $(\mathcal{A}, \text{d})$ be a sheaf of differential graded algebras
on $(\mathcal{C}, \mathcal{O})$. There exists a set $S$ and for each $s$
an acyclic differential graded $\mathcal{A}$-module $\mathcal{M}_s$ such
that for every nonzero acyclic differential graded $\mathcal{A}$-module
$\mathcal{M}$ there is an $s \in S$ and an injective map
$\mathcal{M}_s \to \mathcal{M}$ in $\textit{Mod}(\mathcal{A}, \text{d})$.
\end{lemma}

\begin{proof}
Before we start recall that our conventions guarantee the site $\mathcal{C}$
has a set of objects and morphisms and a set $\text{Cov}(\mathcal{C})$
of coverings.
If $\mathcal{F}$
is a differential graded $\mathcal{A}$-module, let us define $|\mathcal{F}|$
to be the sum of the cardinality of
$$
\coprod\nolimits_{(U, n)} \mathcal{F}^n(U)
$$
as $U$ ranges over the objects of $\mathcal{C}$ and $n \in \mathbf{Z}$.
Choose an infinite cardinal $\kappa$ bigger than
the cardinals $|\Ob(\mathcal{C})|$, $|\text{Arrows}(\mathcal{C})|$,
$|\text{Cov}(\mathcal{C})|$, $\sup |I|$ for
$\{U_i \to U\}_{i \in I} \in \text{Cov}(\mathcal{C})$,
and $|\mathcal{A}|$.

\medskip\noindent
Let $\mathcal{F} \subset \mathcal{M}$ be an inclusion of
differential graded $\mathcal{A}$-modules.
Suppose given a set $K$ and for each $k \in K$ a triple
$(U_k, n_k, x_k)$ consisting of an object $U_k$ of $\mathcal{C}$,
integer $n_k$, and a section $x_k \in \mathcal{M}^{n_k}(U_k)$.
Then we can consider the smallest differential graded
$\mathcal{A}$-submodule $\mathcal{F}' \subset \mathcal{M}$
containing $\mathcal{F}$ and the sections $x_k$ for $k \in K$.
We can describe
$$
(\mathcal{F}')^n(U) \subset \mathcal{M}^n(U)
$$
as the set of elements $x \in \mathcal{M}^n(U)$ such that there
exists $\{f_i : U_i \to U\}_{i \in I} \in \text{Cov}(\mathcal{C})$ such that
for each $i \in I$ there is a finite set $T_i$ and morphisms
$g_t : U_i \to U_{k_t}$
$$
f_i^*x = y_i +
\sum\nolimits_{t \in T_i} a_{it}g_t^*x_{k_t} + b_{it}g_t^*\text{d}(x_{k_t})
$$
for some section $y_i \in \mathcal{F}^n(U)$ and
sections $a_{it} \in \mathcal{A}^{n - n_{k_t}}(U_i)$
and $b_{it} \in \mathcal{A}^{n - n_{k_t} - 1}(U_i)$.
(Details omitted; hints: these sections are certainly in $\mathcal{F}'$
and you show conversely that this rule defines a differential
graded $\mathcal{A}$-submodule.)
It follows from this description that
$|\mathcal{F}'| \leq \max(|\mathcal{F}|, |K|, \kappa)$.

\medskip\noindent
Let $\mathcal{M}$ be a nonzero acyclic differential graded
$\mathcal{A}$-module. Then we can find an integer $n$ and a
nonzero section $x$ of $\mathcal{M}^n$ over some object $U$
of $\mathcal{C}$. Let
$$
\mathcal{F}_0 \subset \mathcal{M}
$$
be the smallest differential graded $\mathcal{A}$-submodule
containing $x$. By the previous paragraph we have
$|\mathcal{F}_0| \leq \kappa$. By induction, given
$\mathcal{F}_0, \ldots, \mathcal{F}_n$ define
$\mathcal{F}_{n + 1}$ as follows. Consider the set
$$
L = \{(U, n, x)\}
\{U_i \to U\}_{i \in I}, (x_i)_{i \in I})\}
$$
of triples where $U$ is an object of $\mathcal{C}$, $n \in \mathbf{Z}$,
and $x \in \mathcal{F}_n(U)$ with $\text{d}(x) = 0$. Since
$\mathcal{M}$ is acyclic for each triple $l = (U_l, n_l, x_l) \in L$
we can choose
$\{(U_{l, i} \to U_l\}_{i \in I_l} \in \text{Cov}(\mathcal{C})$ and
$x_{l, i} \in \mathcal{M}^{n_l - 1}(U_{l, i})$ such that
$\text{d}(x_{l, i}) = x|_{U_{l, i}}$. Then we set
$$
K = \{(U_{l, i}, n_l - 1, x_{l, i}) \mid l \in L, i \in I_l\}
$$
and we let $\mathcal{F}_{n + 1}$ be the smallest differential
graded $\mathcal{A}$-submodule of $\mathcal{M}$ containing
$\mathcal{F}_n$ and the sections $x_{l, i}$.
Since $|K| \leq \max(\kappa, |\mathcal{F}_n|)$
we conclude that $|\mathcal{F}_{n + 1}| \leq \kappa$ by induction.

\medskip\noindent
By construction the inclusion $\mathcal{F}_n \to \mathcal{F}_{n + 1}$
induces the zero map on cohomology sheaves. Hence we see that
$\mathcal{F} = \bigcup \mathcal{F}_n$ is a nonzero acyclic submodule
with $|\mathcal{F}| \leq \kappa$. Since there is only a set
of isomorphism classes of differential graded $\mathcal{A}$-modules
$\mathcal{F}$ with $|\mathcal{F}|$ bounded, we conclude.
\end{proof}

\end{document}
